% ============================================================
% 第一次测试 6 题解题报告
%
% 编译（中文必须用 xelatex，且编两遍目录才准）：
%     cd ~/tjurm2027-test1/report
%     xelatex main.tex && xelatex main.tex
%     （或 latexmk -xelatex -pdf main.tex）
%
% 填写纪律：每节固定四小节 —— 题目 / 思路 / 代码 / 结果。
%   题目 = 用自己的话复述（别看源码）
%   思路 = 一两句话讲清算法怎么走
%   代码 = 粘 src/tests.cpp 里的实现（或用 \lstinputlisting 引用）
%   结果 = 测试输出 / 结果图
% 把下面所有 “TODO” 换成你自己的内容即可。
% ============================================================
\documentclass[11pt,a4paper]{ctexart}

\usepackage[margin=2.5cm]{geometry}
\usepackage{graphicx}
\usepackage{listings}
\usepackage{amsmath}
\usepackage{booktabs}
\usepackage{xcolor}
\usepackage{hyperref}
\hypersetup{colorlinks=true, linkcolor=blue, urlcolor=blue}

\lstset{
  basicstyle=\ttfamily\small,
  numbers=left,
  numberstyle=\tiny\color{gray},
  frame=single,
  breaklines=true,
  language=C++,
  keywordstyle=\color{blue}\bfseries,
  commentstyle=\color{teal},
  stringstyle=\color{orange!80!black},
  showstringspaces=false,
  columns=flexible,
}

\title{第一次测试 6 题解题报告}
\author{Liboven \\ \small 北洋机甲算法组 2027 赛季培训}
\date{\today}

\begin{document}

% 目录只列到 section 级（6 题 + 附录），否则 30 行小节会把目录挤到第二页
\setcounter{tocdepth}{1}

\maketitle

\begin{abstract}
本报告记录我对北洋机甲算法组第一次测试 6 道题的实现思路与验证结果：前三题用 C++ 复现
C 字符串库函数（\texttt{my\_strlen}、\texttt{my\_strcat}、\texttt{my\_strstr}），
后三题是图像处理中的常用函数（\texttt{rgb2gray}、\texttt{resize}、\texttt{hist\_eq}）。
所有代码在 WSL Ubuntu 22.04 下用 CMake + g++ 11.4 编译；图像结果由程序生成，
并与仓库中的参考答案逐像素比对。
\end{abstract}

\tableofcontents
\newpage

% ============================================================
\section{my\_strlen：自己实现字符串长度}
\subsection{题目}
% TODO：用自己的话复述要求（别看源码，凭理解写）
统计一个以 \texttt{'\textbackslash 0'} 结尾的字符串有多少个字符，不调用现成的
\texttt{strlen}。

\subsection{思路}
% TODO：一两句话。指针怎么走、什么时候停
从首地址开始，只要当前字符不是结束符就往后走一步并计数；遇到结束符时计数值就是长度。

\subsection{代码}
\begin{lstlisting}
int my_strlen(char* str) {
  int len = 0;
  while (str[len] != '\0') ++len;
  return len;
}
\end{lstlisting}

\subsection{结果}
% TODO：换成你的终端输出
\texttt{test\_strlen()} 输出：\textbf{通过}（对 \texttt{"123456"}、\texttt{""}、
\texttt{"hello world!"} 三个样例都与标准库结果一致）。

% ============================================================
\section{my\_strcat：字符串拼接}
\subsection{题目}
% TODO
\subsection{思路}
% TODO
\subsection{代码}
\begin{lstlisting}
// TODO：粘上你的 my_strcat 实现
\end{lstlisting}
\subsection{结果}
% TODO

% ============================================================
\section{my\_strstr：子串查找}
\subsection{题目}
% TODO
\subsection{思路}
% TODO
\subsection{代码}
\begin{lstlisting}
// TODO：粘上你的 my_strstr 实现
\end{lstlisting}
\subsection{结果}
% TODO

% ============================================================
\section{rgb2gray：彩色图转灰度图}
\subsection{题目}
% TODO
\subsection{思路}
% TODO：记得写清下标公式
彩色图在内存里按 \texttt{R,G,B} 三通道连续存放，第 $(x,y)$ 个像素的三个通道位于
\texttt{in[(y*w+x)*3]}、\texttt{+1}、\texttt{+2}。灰度公式为
\begin{equation}
V = 0.2989\,R + 0.5870\,G + 0.1140\,B
\end{equation}

\subsection{代码}
\begin{lstlisting}
// TODO：粘上你的 rgb2gray 实现
\end{lstlisting}

\subsection{结果}
\begin{figure}[h]
  \centering
  \includegraphics[width=0.45\textwidth]{../images/rgb2gray/input.jpg}
  \includegraphics[width=0.45\textwidth]{../images/rgb2gray/output.jpg}
  \caption{左：输入彩色图；右：我的程序输出的灰度图}
\end{figure}

% ============================================================
\section{resize：双线性插值缩放}
\subsection{题目}
% TODO
\subsection{思路}
% TODO
\subsection{代码}
\begin{lstlisting}
// TODO：粘上你的 resize 实现
\end{lstlisting}
\subsection{结果}
% TODO：可插入 ../images/resize/output0.jpg 与 output1.jpg

% ============================================================
\section{hist\_eq：直方图均衡化}
\subsection{题目}
% TODO
\subsection{思路}
% TODO
\subsection{代码}
\begin{lstlisting}
// TODO：粘上你的 hist_eq 实现
\end{lstlisting}
\subsection{结果}
% TODO

% ============================================================
\section*{附录：如何验证}
\addcontentsline{toc}{section}{附录：如何验证}
\begin{lstlisting}[language=bash]
cd ~/tjurm2027-test1
cmake --build build -j$(nproc)
cd build && ./tjurm-test          # 注意必须在 build/ 里运行（图片路径是 ../ 相对的）
\end{lstlisting}
图像题的结果可用 \texttt{git status} 做字节级验收：产出图若与仓库里的参考版本不同，
它们会显示为修改状态。

\end{document}
